% GENERATED FILE. Edit canonical lessons, then run npm run generate:books.
% canonical-source-hash: fnv1a64:c6b8862ee3ef9b3a
% canonical-lessons: GE-C401-hupe, GE-C401-hupen, GE-C401-sicherheitsgurt, GE-C401-anschnallen, GE-C401-beifahrer

\chapter{Sicher unterwegs}
\label{ch:ge-sicher-unterwegs}
\hlchaptermodality{401}

\begin{chapteropening}
\textbf{By the end of this chapter:} \emph{I can talk about safety in the car: the horn, hooting, the seat belt, buckling up and the front-seat passenger.}

Complete the last lesson of chapter 401: I can talk about safety in the car: the horn, hooting, the seat belt, buckling up and the front-seat passenger.
\end{chapteropening}

\section[die Hupe]{die Hupe — the horn (of a car)}

\label{lesson:GE-C401-hupe}



\begin{quote}
Before the new one: say the German for the rear-view mirror, then the German for to indicate.
\end{quote}

\subsection*{You'll want to know: die Hupe}
\textbf{die Hupe} — \textquotedblleft{}the horn (of a car)\textquotedblright{}.

\begin{quote}
Die Hupe von unserem alten Auto ist sehr laut.
\end{quote}

\begin{quote}
\emph{The horn on our old car is very loud.}
\end{quote}

Feminine: \emph{eine laute Hupe}.

\begin{cousinweb}
The origin of \emph{Hupe} is not settled. It is usually explained as an imitation of the sound itself; EtymDB does not list the word.
\end{cousinweb}

\begin{grammarlens}[title={What you've built}]
One more word for telling, explaining and describing at B1.
\end{grammarlens}

\subsection*{Your turn}
\begin{itemize}
\raggedright
  \item \emph{Say it:} \emph{die Hupe}
  \item \emph{Say it:} \emph{die Hupe}, once more
  \item \emph{Say it:} a sentence of your own with \emph{die Hupe}
  \item \emph{Recall:} say the German for the rear-view mirror, then the German for to indicate, then say \emph{die Hupe} again
\end{itemize}

\begin{tcolorbox}[breakable,colback=teal!4,colframe=teal!35!black,title={Before you move on}]
What does die Hupe mean? (\textquotedblleft{}The horn (of a car)\textquotedblright{}.) Say it once more.
\end{tcolorbox}
\section[hupen]{hupen — to hoot, to sound the horn}

\label{lesson:GE-C401-hupen}



\begin{quote}
Before the new one: say the German for to indicate, then the German for the horn.
\end{quote}

\subsection*{You'll want to know: hupen}
\textbf{hupen} — \textquotedblleft{}to hoot, to sound the horn\textquotedblright{}.

\begin{quote}
Der Taxifahrer hupt zweimal vor unserem Haus.
\end{quote}

\begin{quote}
\emph{The taxi driver hoots twice in front of our house.}
\end{quote}

Regular: \emph{er hupt}, \emph{er hat gehupt}.

\begin{cousinweb}
\emph{hupen} is built on \emph{Hupe} 'horn', which is usually explained as a sound word. EtymDB lists neither, so the origin stays hedged.
\end{cousinweb}

\begin{grammarlens}[title={What you've built}]
One more word for telling, explaining and describing at B1.
\end{grammarlens}

\subsection*{Your turn}
\begin{itemize}
\raggedright
  \item \emph{Say it:} \emph{hupen}
  \item \emph{Say it:} \emph{hupen}, once more
  \item \emph{Say it:} a sentence of your own with \emph{hupen}
  \item \emph{Recall:} say the German for to indicate, then the German for the horn, then say \emph{hupen} again
\end{itemize}

\begin{tcolorbox}[breakable,colback=teal!4,colframe=teal!35!black,title={Before you move on}]
What does hupen mean? (\textquotedblleft{}To hoot, to sound the horn\textquotedblright{}.) Say it once more.
\end{tcolorbox}
\section[der Sicherheitsgurt]{der Sicherheitsgurt — the seat belt}

\label{lesson:GE-C401-sicherheitsgurt}



\begin{quote}
Before the new one: say the German for the horn, then the German for to hoot.
\end{quote}

\subsection*{You'll want to know: der Sicherheitsgurt}
\textbf{der Sicherheitsgurt} — \textquotedblleft{}the seat belt\textquotedblright{}.

\begin{quote}
Im Taxi trägt sie immer den Sicherheitsgurt.
\end{quote}

\begin{quote}
\emph{In a taxi she always wears her seat belt.}
\end{quote}

Plural: \emph{die Sicherheitsgurte}. Often shortened to \emph{der Gurt}.

\begin{cousinweb}
\emph{Sicherheitsgurt} is built from \emph{Sicherheit} 'safety', from \emph{sicher}, an early borrowing from Latin \emph{sēcūrus} 'free from care' (Lewis \& Short), and \emph{Gurt} 'belt', inherited from Middle High German \emph{gurt} (EtymDB).
\end{cousinweb}

\begin{grammarlens}[title={What you've built}]
One more word for telling, explaining and describing at B1.
\end{grammarlens}

\subsection*{Your turn}
\begin{itemize}
\raggedright
  \item \emph{Say it:} \emph{der Sicherheitsgurt}
  \item \emph{Say it:} \emph{der Sicherheitsgurt}, once more
  \item \emph{Say it:} a sentence of your own with \emph{der Sicherheitsgurt}
  \item \emph{Recall:} say the German for the horn, then the German for to hoot, then say \emph{der Sicherheitsgurt} again
\end{itemize}

\begin{tcolorbox}[breakable,colback=teal!4,colframe=teal!35!black,title={Before you move on}]
What does der Sicherheitsgurt mean? (\textquotedblleft{}The seat belt\textquotedblright{}.) Say it once more.
\end{tcolorbox}
\section[anschnallen]{anschnallen — to buckle up, to fasten a seat belt (mostly reflexive: sich anschnallen)}

\label{lesson:GE-C401-anschnallen}



\begin{quote}
Before the new one: say the German for to hoot, then the German for the seat belt.
\end{quote}

\subsection*{You'll want to know: anschnallen}
\textbf{anschnallen} — \textquotedblleft{}to buckle up, to fasten a seat belt (mostly reflexive: sich anschnallen)\textquotedblright{}.

\begin{quote}
Vor der Fahrt schnallen sich alle Kinder an.
\end{quote}

\begin{quote}
\emph{Before the journey all the children buckle up.}
\end{quote}

Mostly reflexive and separable: \emph{ich schnalle mich an}, \emph{wir haben uns angeschnallt}.

\begin{cousinweb}
\emph{anschnallen} is built from \emph{an} 'on' and \emph{Schnalle} 'buckle', which is inherited from Middle High German \emph{snalle} (EtymDB).
\end{cousinweb}

\begin{grammarlens}[title={What you've built}]
One more word for telling, explaining and describing at B1.
\end{grammarlens}

\subsection*{Your turn}
\begin{itemize}
\raggedright
  \item \emph{Say it:} \emph{anschnallen}
  \item \emph{Say it:} \emph{anschnallen}, once more
  \item \emph{Say it:} a sentence of your own with \emph{anschnallen}
  \item \emph{Recall:} say the German for to hoot, then the German for the seat belt, then say \emph{anschnallen} again
\end{itemize}

\begin{tcolorbox}[breakable,colback=teal!4,colframe=teal!35!black,title={Before you move on}]
What does anschnallen mean? (\textquotedblleft{}To buckle up, to fasten a seat belt (mostly reflexive: sich anschnallen)\textquotedblright{}.) Say it once more.
\end{tcolorbox}
\section[der Beifahrer]{der Beifahrer — the front-seat passenger}

\label{lesson:GE-C401-beifahrer}



\begin{quote}
Before the new one: say the German for the seat belt, then the German for to buckle up.
\end{quote}

\subsection*{You'll want to know: der Beifahrer}
\textbf{der Beifahrer} — \textquotedblleft{}the front-seat passenger\textquotedblright{}.

\begin{quote}
Der Beifahrer liest die Karte und sagt den Weg.
\end{quote}

\begin{quote}
\emph{The passenger reads the map and says which way to go.}
\end{quote}

Same form in the plural. A woman: \emph{die Beifahrerin}. \emph{Der Beifahrersitz} — the passenger seat.

\begin{cousinweb}
\emph{Beifahrer} is built from \emph{bei} 'beside, with' and \emph{Fahrer} 'driver', from the inherited verb \emph{fahren}, Proto-Germanic \emph{faraną} (EtymDB).
\end{cousinweb}

\begin{grammarlens}[title={What you've built}]
One more word for telling, explaining and describing at B1.
\end{grammarlens}

\subsection*{Your turn}
\begin{itemize}
\raggedright
  \item \emph{Say it:} \emph{der Beifahrer}
  \item \emph{Say it:} \emph{der Beifahrer}, once more
  \item \emph{Say it:} a sentence of your own with \emph{der Beifahrer}
  \item \emph{Recall:} say the German for the seat belt, then the German for to buckle up, then say \emph{der Beifahrer} again
\end{itemize}

\begin{tcolorbox}[breakable,colback=teal!4,colframe=teal!35!black,title={Before you move on}]
What does der Beifahrer mean? (\textquotedblleft{}The front-seat passenger\textquotedblright{}.) Say it once more.
\end{tcolorbox}
